% ============================================================
%  CHAPITRE 2 : De la Donnée Brute à la Donnée Prête
% ============================================================
\chapter{De la Donnée Brute à la Donnée Prête}
\label{chap:data}

\begin{infobox}[Résumé du chapitre]
Ce chapitre détaille la transformation des fichiers CMS bruts en un dataset propre,
structuré et exempt de fuites de données. Il couvre l'analyse exploratoire,
le nettoyage, la construction anti-leakage de la variable cible, le split temporel,
l'ingénierie des 29 variables finales, la validation Great Expectations et le
versionnement DVC.
\end{infobox}

% ============================================================
\section{Analyse exploratoire des données (EDA)}
% ============================================================

\subsection{Vue d'ensemble statistique}

L'analyse exploratoire a porté sur la table des bénéficiaires fusionnée sur trois
années (2008--2010), représentant 88\,585~lignes et 33~colonnes issues du croisement
des fichiers patients, hospitalisations, consultations, actes médicaux et prescriptions.
Cette étape a permis de qualifier la richesse des données disponibles, d'identifier
les anomalies à corriger et de comprendre la structure temporelle du dataset avant
toute modélisation.

\begin{table}[htbp]
\centering
\caption{Statistiques descriptives de la cohorte}
\label{tab:eda_stats}
\begin{tabular}{@{}ll@{}}
\toprule
\textbf{Indicateur} & \textbf{Valeur} \\
\midrule
Patients uniques         & 30\,000 \\
Doublons                 & 0 \\
Lignes fusionnées        & 88\,585 (3~années $\times$ patients) \\
Colonnes                 & 33 \\
Femmes / Hommes          & 55.3\,\% / 44.7\,\% \\
Décès enregistrés        & 490 \\
Plage de dates           & 2007-12-10 à 2010-12-30 \\
Taux d'hospitalisation   & 8.10\,\% global \\
\bottomrule
\end{tabular}
\end{table}

La répartition par genre est légèrement déséquilibrée au profit des femmes, ce
qui reflète la démographie réelle de la population Medicare. L'absence totale de
doublons sur la clé patient-année confirme la cohérence du schéma de fusion
multi-annuelle.


\subsection{Taux d'hospitalisation par année}

Une analyse temporelle révèle une variation non négligeable du taux d'hospitalisation
entre les trois années. En 2009, le taux grimpe à 9.18\,\%, probablement sous l'effet
d'un vieillissement progressif de la cohorte et d'une accumulation de comorbidités
non traitées. Cette hétérogénéité temporelle justifie un split temporel strict plutôt
qu'une validation croisée classique, qui mélangerait les années et biaiserait
l'évaluation des performances réelles du modèle.

\begin{table}[htbp]
\centering
\caption{Taux d'hospitalisation par année}
\label{tab:hosp_annee}
\begin{tabular}{@{}llll@{}}
\toprule
\textbf{Année} & \textbf{Patients} & \textbf{Hospitalisés} & \textbf{Taux} \\
\midrule
2008 (Train)       & 30\,000 & 2\,261 & 7.54\,\% \\
2009 (Validation)  & 29\,510 & 2\,710 & 9.18\,\% \\
2010 (Test)        & 29\,075 & 2\,203 & 7.58\,\% \\
\midrule
\textbf{Global}    & 88\,585 & 7\,174 & \textbf{8.10\,\%} \\
\bottomrule
\end{tabular}
\end{table}

\subsection{Distribution des comorbidités chroniques}

Les 11~comorbidités chroniques codées dans les données présentent des prévalences
très variables, de 2.1\,\% pour la maladie rénale chronique jusqu'à 43.8\,\% pour
la maladie d'Alzheimer. Ce gradient de prévalence est cliniquement cohérent~: les
pathologies neurodégénératives et cardiovasculaires dominent chez la population
Medicare, âgée en moyenne de 70 ans ou plus. La relation entre comorbidités et
hospitalisation est quantifiée par l'odds ratio, qui mesure combien de fois le
risque d'hospitalisation est multiplié en présence de la pathologie. L'insuffisance
cardiaque présente l'odds ratio le plus élevé (5.1), ce qui en fait le prédicteur
clinique le plus fort, devant la BPCO (3.1) et l'AVC/AIT (3.0).

\begin{table}[htbp]
\centering
\caption{Comorbidités chroniques~: prévalence et association à l'hospitalisation}
\label{tab:comorbidites}
\begin{tabular}{@{}llrrl@{}}
\toprule
\textbf{Code} & \textbf{Maladie} & \textbf{Prévalence} & \textbf{Poids} & \textbf{Odds Ratio} \\
\midrule
SP\_CHF       & Insuf. cardiaque  & 41.3\,\% & 1 & 5.1 \\
SP\_ALZHDMTA  & Alzheimer         & 43.8\,\% & 1 & --- \\
SP\_RA\_OA    & Arthrite          & 38.5\,\% & 1 & --- \\
SP\_ISCHMCHT  & Coronaropathie    & 31.2\,\% & 1 & 2.5 \\
SP\_DIABETES  & Diabète           & 28.7\,\% & 1 & 2.3 \\
SP\_OSTEOPRS  & Ostéoporose       & 22.4\,\% & 0 & --- \\
SP\_COPD      & BPCO              & 18.9\,\% & 1 & 3.1 \\
SP\_DEPRESSN  & Dépression        & 14.2\,\% & 1 & 1.9 \\
SP\_STRKETIA  & AVC/AIT           &  8.3\,\% & 2 & 3.0 \\
SP\_CNCR      & Cancer            &  5.8\,\% & 2 & --- \\
SP\_CHRNKIDN  & Maladie rénale    &  2.1\,\% & 2 & --- \\
\bottomrule
\end{tabular}
\end{table}

\subsection{Résultats ACP et segmentation KMeans}

Une analyse en composantes principales (ACP) et une segmentation KMeans en
3~clusters ont permis d'identifier des profils cliniques distincts au sein de la
cohorte. Les trois premières composantes principales expliquent 62.1\,\% de la
variance totale (CP1~: 31.4\,\%, CP2~: 18.7\,\%), ce qui indique une structure
latente significative dans les données. La segmentation révèle trois niveaux de
risque naturellement séparables, validant l'hypothèse que les patients
polypathologiques forment une population hétérogène stratifiable.

\begin{resultatbox}[Segmentation des patients en 3 clusters]
\begin{tabular}{@{}llll@{}}
\toprule
\textbf{Cluster} & \textbf{Patients} & \textbf{Taux hosp.} & \textbf{Comorbidités moy.} \\
\midrule
0 — Faible risque   & 14\,820 & 3.2\,\%  & 0.8 \\
1 — Risque modéré   & 10\,940 & 7.8\,\%  & 2.3 \\
2 — Haut risque     &  4\,240 & 22.1\,\% & 5.1 \\
\bottomrule
\end{tabular}

\vspace{0.5em}
L'ACP explique 62.1\,\% de la variance totale sur 3~composantes
(CP1~: 31.4\,\%, CP2~: 18.7\,\%).
\end{resultatbox}

Le cluster à haut risque (4\,240 patients, 22.1\,\% de taux d'hospitalisation)
représente la population cible prioritaire des interventions préventives. Ces patients
cumulent en moyenne 5.1~comorbidités actives, ce qui correspond au profil typique du
patient polypathologique décrit dans la littérature clinique.


% ------------------------------------------------------------------
%  FIGURE 2.1 — Distribution + taux hospitalisation par âge
% ------------------------------------------------------------------
\subsection{Distribution de la cible et taux par groupe d'âge}

\begin{figure}[htbp]
\centering
\begin{subfigure}[b]{0.48\textwidth}
  \begin{tikzpicture}
    \begin{axis}[
      ybar, bar width=20pt,
      xlabel={Statut}, ylabel={Nombre de patients},
      symbolic x coords={Non hospitalisé, Hospitalisé},
      xtick=data, ymin=0, ymax=85000,
      nodes near coords,
      nodes near coords align={vertical},
      every node near coord/.append style={font=\tiny},
      width=\textwidth, height=5.5cm,
      title={Distribution de la variable cible},
      title style={font=\small},
      tick label style={font=\small},
      label style={font=\small},
    ]
    \addplot[fill=primary!70] coordinates {
      (Non hospitalisé, 81411)
      (Hospitalisé, 7174)
    };
    \end{axis}
  \end{tikzpicture}
  \caption{Déséquilibre 91.9\,\% / 8.1\,\%}
\end{subfigure}
\hfill
\begin{subfigure}[b]{0.48\textwidth}
  \begin{tikzpicture}
    \begin{axis}[
      ybar, bar width=12pt,
      xlabel={Groupe d'âge},
      ylabel={Taux d'hospitalisation (\%)},
      symbolic x coords={<65, 65-74, 75-84, 85+},
      xtick=data, ymin=0, ymax=22,
      nodes near coords,
      nodes near coords align={vertical},
      every node near coord/.append style={font=\tiny},
      width=\textwidth, height=5.5cm,
      title={Taux hosp. par groupe d'âge},
      title style={font=\small},
      tick label style={font=\small},
      label style={font=\small},
    ]
    \addplot[fill=accent!80] coordinates {
      (<65,3.8)(65-74,6.2)(75-84,11.4)(85+,19.7)
    };
    \end{axis}
  \end{tikzpicture}
  \caption{Croissance du risque avec l'âge}
\end{subfigure}
\caption{Distribution de la variable cible (a) et taux d'hospitalisation par groupe d'âge (b)}
\label{fig:distrib_cible}
\end{figure}

L'histogramme de la figure~\ref{fig:distrib_cible}(a) confirme le déséquilibre
marqué de la variable cible~: 91.9\,\% de patients non hospitalisés contre 8.1\,\%.
Ce ratio de 1:12 constitue l'un des défis centraux de la modélisation et justifie
l'utilisation de stratégies de pondération des classes. La progression exponentielle
du taux d'hospitalisation avec l'âge~--- de 3.8\,\% avant 65~ans à 19.7\,\% après
85~ans~--- souligne l'importance de l'âge comme facteur de risque clinique fondamental.

% ============================================================
\section{Nettoyage des données}
% ============================================================

\subsection{Problèmes identifiés et solutions appliquées}

L'analyse exploratoire a mis en évidence quatre catégories de problèmes de qualité
dans les données brutes CMS. Ces anomalies, si elles n'étaient pas corrigées, auraient
introduit des biais silencieux dans les modèles ou causé des erreurs lors du traitement
du pipeline. Chaque problème a été traité par une stratégie ciblée, documentée et
reproductible.

Le premier problème concernait le format des dates~: certaines colonnes de dates
(naissance, décès, hospitalisation) étaient stockées sous forme de nombres flottants
du type \texttt{20100312.0} au lieu de chaînes de caractères. L'interprétation directe
de ces valeurs par les bibliothèques pandas produisait des timestamps incohérents
(remontant à 1970), corrompant silencieusement les calculs de fenêtres temporelles.
La solution retenue consiste à isoler la partie entière de la valeur flottante, puis
à la parser selon le format \texttt{\%Y\%m\%d}. L'implémentation de cette correction,
ainsi que les commandes DVC associées, est disponible en Annexe~A.

Le deuxième problème portait sur la colonne \texttt{BENE\_ESRD\_IND} (indicateur
d'insuffisance rénale terminale)~: les valeurs mélangent des chaînes de caractères
\texttt{'Y'} et des entiers \texttt{0}, rendant la colonne inutilisable directement
comme feature numérique. Un mapping explicite \texttt{Y}~$\rightarrow$~1, reste~$\rightarrow$~0
a standardisé la colonne en variable binaire cohérente.

Le troisième problème concernait les valeurs extrêmes dans les neuf colonnes de coûts
médicaux. Certains patients présentaient des remboursements anormalement élevés, probablement
liés à des hospitalisations chirurgicales lourdes ou à des traitements onéreux. Ces outliers
auraient déstabilisé l'entraînement des modèles par gradient boosting. La winsorisation au
percentile~99 limite l'influence de ces valeurs sans les supprimer, en les remplaçant par
la valeur maximale acceptable. Le code correspondant est fourni en Annexe~A.

Enfin, la fusion des données sur trois années a créé des doublons pour les patients présents
dans plusieurs fichiers annuels. Une clé composite \texttt{(DESYNPUF\_ID, ANNEE)} a permis
d'identifier et de supprimer ces doublons tout en préservant l'historique multi-annuel.

\begin{table}[htbp]
\centering
\caption{Problèmes identifiés lors de l'EDA et solutions appliquées}
\label{tab:bugs_eda}
\begin{tabular}{@{}p{3cm}p{4.5cm}p{5cm}@{}}
\toprule
\textbf{Problème} & \textbf{Manifestation} & \textbf{Solution} \\
\midrule
Dates flottants  & \texttt{20100312.0} → timestamp 1970
                 & Extraction partie entière + format \texttt{\%Y\%m\%d} \\
BENE\_ESRD\_IND  & Valeurs 'Y' et '0' mélangées
                 & Mapping \texttt{Y}→1, autres→0 \\
Coûts extrêmes   & Valeurs aberrantes ($>$P99)
                 & Winsorisation au percentile~99 \\
Doublons multi-années & Même patient dans 2009 et 2010
                 & Clé unique \texttt{(DESYNPUF\_ID, ANNEE)} \\
\bottomrule
\end{tabular}
\end{table}

\subsection{Résultats du nettoyage}

À l'issue de cette phase, le dataset nettoyé présente les caractéristiques suivantes.

\begin{resultatbox}[Table finale après nettoyage]
\begin{itemize}
  \item 88\,585~lignes, 33~colonnes
  \item Aucun doublon sur la clé \texttt{(DESYNPUF\_ID, ANNEE)}
  \item 6\,216~patients cold-start (\texttt{IS\_NEW\_PATIENT}~=~1)
  \item 9~colonnes de coûts winsorisées au P99
\end{itemize}
\end{resultatbox}

Les 6\,216 patients cold-start correspondent à des bénéficiaires n'ayant aucun
historique médical dans les données disponibles. Plutôt que de les exclure, un
indicateur binaire \texttt{IS\_NEW\_PATIENT} a été créé pour permettre au modèle
d'apprendre le comportement spécifique de cette sous-population, généralement
moins bien servie par les systèmes de prédiction traditionnels.


% ============================================================
\section{Construction anti-leakage de la variable cible}
% ============================================================

\subsection{Le risque de fuite temporelle}

La construction naïve de la variable cible constitue l'un des pièges les plus courants
du machine learning médical. Si l'on définit simplement~: \og le patient a-t-il été
hospitalisé pendant l'année $t$~?\fg{} en utilisant des features calculées sur cette
même année, le modèle peut exploiter des informations sur des hospitalisations survenues
\textit{après} la période d'observation. Cette fuite de données temporelle
(\textit{data leakage}) produit des modèles avec des métriques artificiellement gonflées,
qui se révèlent inefficaces en déploiement réel.

\subsection{Définition formelle de la cible}

Pour éliminer ce risque, la variable cible est définie avec une séparation stricte
entre fenêtre d'observation et fenêtre de prédiction~:

\begin{equation}
\label{eq:cible}
y_i =
\begin{cases}
1 & \text{si le patient } i \text{ est hospitalisé entre } t+1 \text{ et } t+365 \\
0 & \text{sinon}
\end{cases}
\end{equation}

\noindent Les features sont calculées exclusivement sur la fenêtre $[t-365, t]$, et
la cible sur la fenêtre strictement postérieure $[t+1, t+365]$. Il n'y a aucun
chevauchement entre les deux fenêtres. Concrètement, pour l'entraînement sur 2008~:
les features décrivent le comportement médical du patient sur l'année 2007--2008,
et la cible indique s'il a été hospitalisé au cours de l'année 2009.

\subsection{Fenêtres temporelles anti-leakage}

La figure~\ref{fig:fenetre_temporelle} illustre le principe des fenêtres temporelles
appliquées aux trois splits du dataset.

\begin{figure}[htbp]
\centering
% ------------------------------------------------------------------
%  FIGURE 2.2 — Fenêtres temporelles (TikZ)
% ------------------------------------------------------------------
\begin{tikzpicture}[font=\small, >=Stealth]

% Axe du temps
\draw[thick, ->] (-0.5,0) -- (12,0) node[right] {Temps};
\foreach \x/\l in {0/{jan. 2007}, 3/{jan. 2008}, 6/{jan. 2009}, 9/{jan. 2010}, 11.5/{déc. 2010}}
  \draw (\x, 0.1) -- (\x,-0.1) node[below, font=\tiny] {\l};

% Train 2008
\fill[primary!30] (0,0.5) rectangle (3,1.5);
\fill[success!30]  (3,0.5) rectangle (6,1.5);
\draw[thick, primary] (0,0.5) rectangle (3,1.5);
\draw[thick, success]  (3,0.5) rectangle (6,1.5);
\node[font=\footnotesize] at (1.5,1.0) {Features};
\node[font=\footnotesize] at (4.5,1.0) {Cible $y$};
\node[left, font=\footnotesize\bfseries] at (0,1.0) {Train 2008};

% Val. 2009
\fill[primary!30] (3,2.0) rectangle (6,3.0);
\fill[success!30]  (6,2.0) rectangle (9,3.0);
\draw[thick, primary] (3,2.0) rectangle (6,3.0);
\draw[thick, success]  (6,2.0) rectangle (9,3.0);
\node[font=\footnotesize] at (4.5,2.5) {Features};
\node[font=\footnotesize] at (7.5,2.5) {Cible $y$};
\node[left, font=\footnotesize\bfseries] at (3,2.5) {Val. 2009};

% Test 2010
\fill[primary!30] (6,3.5) rectangle (9,4.5);
\fill[success!30]  (9,3.5) rectangle (11.5,4.5);
\draw[thick, primary] (6,3.5) rectangle (9,4.5);
\draw[thick, success]  (9,3.5) rectangle (11.5,4.5);
\node[font=\footnotesize] at (7.5,4.0) {Features};
\node[font=\footnotesize] at (10.25,4.0) {Cible $y$};
\node[left, font=\footnotesize\bfseries] at (6,4.0) {Test 2010};

% Légende
\fill[primary!30] (0,-1.2) rectangle (0.6,-0.7);
\node[right] at (0.7,-0.95) {\footnotesize Fenêtre observation (features)};
\fill[success!30] (5,-1.2) rectangle (5.6,-0.7);
\node[right] at (5.7,-0.95) {\footnotesize Fenêtre cible (hospitalisation)};

\end{tikzpicture}
\caption{Fenêtres temporelles anti-leakage~: features et cible ne se chevauchent jamais}
\label{fig:fenetre_temporelle}
\end{figure}

% ============================================================
\section{Split temporel strict}
% ============================================================

\begin{alertbox}[Pourquoi pas de K-Fold Cross-Validation~?]
Dans un contexte de prédiction temporelle, la validation croisée classique
crée une fuite~: les données futures informent l'entraînement passé. Le split
temporel strict garantit que le modèle n'a \textit{jamais vu} les données
futures lors de l'entraînement~:
\begin{itemize}
  \item \textbf{Train~:} 2008 (30\,000 patients, 7.54\,\% hospitalisés)
  \item \textbf{Validation~:} 2009 (29\,510 patients, 9.18\,\%)
  \item \textbf{Test~:} 2010 (29\,075 patients, 7.58\,\%)
\end{itemize}
\end{alertbox}

Ce choix entraîne une conséquence importante~: les hyperparamètres sont optimisés
sur l'ensemble de validation 2009, et l'évaluation finale est réalisée exclusivement
sur le jeu de test 2010, qui n'a \textit{jamais été vu} pendant la phase d'optimisation.
Cette rigueur garantit que l'AUC de 0.897 reportée sur le test est une estimation
non biaisée des performances réelles en déploiement.


% ============================================================
\section{Feature Engineering~: les 29 variables finales}
% ============================================================

\subsection{Trois catégories de features}

L'ingénierie des features est l'étape qui transforme les données administratives
brutes en signaux cliniquement interprétables. Les 29~features finales se répartissent
en trois catégories complémentaires, chacune capturant une dimension différente du
profil de risque d'un patient.

\paragraph{Features statiques (16 variables)}
Ces features décrivent les caractéristiques démographiques et cliniques stables d'un
patient~: âge, sexe encodé, groupe racial encodé, indicateur d'insuffisance rénale
terminale (\texttt{BENE\_ESRD\_IND}), groupe d'âge encodé et les 11~indicateurs binaires
de comorbidités chroniques (\texttt{SP\_*}). Bien que \og statiques \fg{} dans le sens
où elles varient peu d'une année à l'autre, elles constituent le profil de risque de base
du patient.

L'index de Charlson (\texttt{CHARLSON\_INDEX}) est calculé à partir de ces comorbidités~:
il attribue un poids de~2 aux pathologies les plus graves (maladie rénale chronique, cancer,
AVC/AIT) et un poids de~1 aux autres. La valeur moyenne dans la cohorte est de 2.34, avec
un maximum observé de~11. La formule de calcul complète est disponible en Annexe~A.

\paragraph{Features dynamiques (8 variables)}
Ces features quantifient l'activité médicale récente d'un patient, calculée par agrégation
sur des fenêtres glissantes de 3, 6 et 12~mois~: nombre d'hospitalisations passées
(\texttt{NB\_HOSP\_PASSEES}), consultations ambulatoires sur 3 mois (\texttt{NB\_OP\_3M}),
6~mois (\texttt{NB\_OP\_6M}) et 12~mois (\texttt{NB\_OP\_12M}), actes médicaux sur 6~mois
(\texttt{NB\_CAR\_6M}), nombre de prescriptions (\texttt{NB\_PRESCRIPTIONS}), nombre de
molécules uniques (\texttt{NB\_MOLECULES\_UNIQUES}) et coût total des hospitalisations
passées (\texttt{COUT\_HOSP\_PASSE}). Ces variables dynamiques captent l'intensité de
consommation de soins, qui s'est révélée être le signal prédictif le plus fort selon
l'analyse SHAP présentée au chapitre suivant.

\paragraph{Features composites (5 variables)}
Ces variables synthétisent plusieurs dimensions en un seul indicateur~:
\texttt{NB\_COMORBIDITES} (compte total des comorbidités actives, moyenne~=~2.26),
\texttt{CHARLSON\_INDEX} (charge de morbidité pondérée, moyenne~=~2.34),
\texttt{COUT\_TOTAL} (agrégation des coûts toutes catégories confondues),
\texttt{POLYPHARMACIE} (indicateur binaire~: $\geq$5~médicaments distincts)
et \texttt{IS\_NEW\_PATIENT} (patient sans historique préalable). Ces features
composites permettent au modèle d'accéder à des résumés de haut niveau sans avoir à
apprendre les interactions depuis zéro.

\begin{table}[htbp]
\centering
\caption{Récapitulatif des 29 features finales}
\label{tab:features}
\begin{tabular}{@{}lll@{}}
\toprule
\textbf{Catégorie} & \textbf{Nb} & \textbf{Variables (exemples)} \\
\midrule
Statiques    & 16 & AGE, SEXE\_ENC, RACE\_ENC, 11 comorbidités SP\_* \\
Dynamiques   &  8 & NB\_HOSP\_PASSEES, NB\_OP\_6M, NB\_PRESCRIPTIONS \\
Composites   &  5 & CHARLSON\_INDEX, NB\_COMORBIDITES, POLYPHARMACIE \\
\midrule
\textbf{Total} & \textbf{29} & \\
\bottomrule
\end{tabular}
\end{table}

% ============================================================
\section{Validation des données~: Great Expectations}
% ============================================================

Great Expectations~\cite{gormley2020great} permet de définir des contrats de données
formels appelés \textit{expectations}, qui s'exécutent automatiquement à chaque passage
dans le pipeline et garantissent la conformité des données avant leur transmission aux
modèles. Trois types de vérifications ont été implémentés~: un contrôle du nombre de
lignes (entre 85\,000 et 95\,000), la présence des colonnes calculées
(\texttt{CHARLSON\_INDEX} et \texttt{NB\_COMORBIDITES}), et un contrôle de cohérence
sur le taux d'hospitalisation moyen (entre 6\,\% et 12\,\%). L'implémentation complète
des \textit{expectations} est disponible en Annexe~A.

\begin{resultatbox}[Résultats Great Expectations]
\begin{itemize}
  \item Test~1 — Nombre de lignes~: \textbf{88\,585} ✓ (attendu~: 85\,000--95\,000)
  \item Test~2 — Colonnes requises~: \textbf{présentes} ✓
  \item Taux de validation~: \textbf{2/2 expectations réussies}
\end{itemize}
\end{resultatbox}

L'intégration de Great Expectations dans le pipeline Prefect garantit que toute
dégradation de qualité des données futures (colonnes manquantes, dérive du taux
d'hospitalisation hors plage, volume insuffisant) sera détectée et bloquée avant
d'atteindre la couche de modélisation.

% ============================================================
\section{Versionnement avec DVC et Git}
% ============================================================

DVC (Data Version Control)~\cite{kupriianov2019dvc} étend Git aux fichiers de données
volumineuses. Le principe est de ne stocker dans Git que de légers fichiers
\texttt{.dvc} contenant le hash MD5 des données, tandis que les données réelles sont
stockées dans un remote (Google Drive, S3, ou stockage local). Chaque version du
dataset est ainsi reproductible à l'identique à partir du hash, garantissant
la traçabilité complète du cycle données~→~modèle~→~prédiction.

Dans ce projet, les répertoires \texttt{data/cleaned/} et \texttt{data/features/} sont
versionnés via DVC, avec un pipeline déclaratif (\texttt{dvc.yaml}) qui rejoue
automatiquement les étapes de nettoyage et de feature engineering lorsque les données
amont changent. Les commandes essentielles sont rassemblées en Annexe~A.

\section*{Synthèse du chapitre}

Ce chapitre a documenté la transformation des données brutes CMS en un dataset
exploitable~: correction de quatre catégories d'anomalies, construction anti-leakage
rigoureuse de la variable cible avec séparation stricte des fenêtres temporelles, et
ingénierie de 29~features cliniquement interprétables. La validation automatisée par
Great Expectations et le versionnement DVC ferment la boucle de qualité~: toute
régression de données est détectée avant d'atteindre les modèles. Le chapitre suivant
explique comment ces données ont été utilisées pour entraîner, optimiser et évaluer
les modèles de prédiction.

